\section{Analysis of the Numerical Results}

\subsection{Overview}

The numerical integration of equation \eqref{eq:scale_factor_ode} yields the evolution of the scale factor $a(t)$ over cosmic time for different values of the decay parameter $\gamma$. This section provides a rigorous mathematical analysis of the computed results.

\begin{definition}[Numerical Solution Space]
Let $\mathcal{S}$ be the space of numerical solutions to the dynamical Friedmann equation:
\[
\mathcal{S} = \{a: [0,T] \to \mathbb{R}^+ \mid a \text{ satisfies } \eqref{eq:scale_factor_ode} \text{ numerically}\},
\]
where $T$ is the integration time horizon.
\end{definition}

\subsection{Key Observations}

\begin{theorem}[Parameter Dependence]
For fixed initial conditions $a(0) = a_0$ and $\dot{a}(0) = H_0 a_0$, the solution $a(t;\gamma)$ satisfies:
\begin{enumerate}
\item Monotonicity: $\frac{\partial a}{\partial \gamma}(t;\gamma) < 0$ for all $t > 0$ and $\gamma \geq 0$.
\item Asymptotic behavior: 
\[
\lim_{t \to \infty} \frac{a(t;\gamma_1)}{a(t;\gamma_2)} = 
\begin{cases}
\infty & \text{if } \gamma_1 < \gamma_2 \\
1 & \text{if } \gamma_1 = \gamma_2 \\
0 & \text{if } \gamma_1 > \gamma_2
\end{cases}
\]
\end{enumerate}
\end{theorem}

\begin{proof}
Consider two solutions $a_1(t) = a(t;\gamma_1)$ and $a_2(t) = a(t;\gamma_2)$ with $\gamma_1 < \gamma_2$. 

For the Hubble parameters:
\[
H_1^2 = 0.3 a_1^{-3} + \Lambda_0 e^{-\gamma_1 t} > 0.3 a_2^{-3} + \Lambda_0 e^{-\gamma_2 t} = H_2^2
\]
when $a_1 = a_2$ (since $e^{-\gamma_1 t} > e^{-\gamma_2 t}$).

Starting from the same initial conditions, $a_1(0) = a_2(0)$, we have $\dot{a}_1(0) > \dot{a}_2(0)$. By continuous dependence on parameters and monotonicity of the flow, $a_1(t) > a_2(t)$ for all $t > 0$.

For the asymptotic behavior, note that as $t \to \infty$:
\[
H_1(t) \approx \sqrt{\Lambda_0 e^{-\gamma_1 t}}, \quad H_2(t) \approx \sqrt{\Lambda_0 e^{-\gamma_2 t}}.
\]
Thus,
\[
\frac{H_1(t)}{H_2(t)} \approx e^{-(\gamma_1-\gamma_2)t/2} \to \infty \text{ as } t \to \infty.
\]
Integrating gives the stated result.
\end{proof}

\subsection{Expansion Dynamics}

\begin{lemma}[Deceleration Parameter]
The deceleration parameter $q(t) = -\frac{\ddot{a}a}{\dot{a}^2}$ for the dynamical dark energy model is given by:
\[
q(t) = \frac{1}{2}\frac{0.3 a^{-3} - 2\Lambda_0 e^{-\gamma t}}{0.3 a^{-3} + \Lambda_0 e^{-\gamma t}}.
\]
\end{lemma}

\begin{proof}
From the Friedmann equation:
\[
H^2 = 0.3 a^{-3} + \Lambda_0 e^{-\gamma t}.
\]
Differentiating with respect to time:
\[
2H\dot{H} = -0.9 a^{-4}\dot{a} - \gamma \Lambda_0 e^{-\gamma t}.
\]
Since $\dot{a} = aH$:
\[
2H\dot{H} = -0.9 a^{-3}H - \gamma \Lambda_0 e^{-\gamma t}.
\]
The acceleration equation:
\[
\frac{\ddot{a}}{a} = \dot{H} + H^2 = -\frac{0.9 a^{-3}}{2} - \frac{\gamma \Lambda_0 e^{-\gamma t}}{2H} + H^2.
\]
After algebraic manipulation:
\[
q = -\frac{\ddot{a}a}{\dot{a}^2} = -\frac{\ddot{a}}{aH^2} = \frac{1}{2}\frac{0.3 a^{-3} - 2\Lambda_0 e^{-\gamma t}}{H^2}.
\]
\end{proof}

\subsection{Impact on Cosmic Expansion}

\begin{proposition}[Transition from Deceleration to Acceleration]
There exists a critical time $t_c$ such that:
\begin{itemize}
\item For $t < t_c$: $q(t) > 0$ (deceleration, matter-dominated)
\item For $t > t_c$: $q(t) < 0$ (acceleration, dark energy-dominated)
\end{itemize}
The critical time is determined by:
\[
0.3 a(t_c)^{-3} = 2\Lambda_0 e^{-\gamma t_c}.
\]
\end{proposition}

\begin{proof}
The transition occurs when $q(t_c) = 0$, which implies:
\[
0.3 a(t_c)^{-3} - 2\Lambda_0 e^{-\gamma t_c} = 0.
\]
Rearranging:
\[
a(t_c)^3 = \frac{0.3}{2\Lambda_0 e^{-\gamma t_c}} = \frac{0.15}{\Lambda_0} e^{\gamma t_c}.
\]
This implicitly defines $t_c$ as a function of the parameters. The existence of such a time is guaranteed by the intermediate value theorem, as $q(0) > 0$ (matter-dominated early universe) and $\lim_{t \to \infty} q(t) < 0$ for $\gamma < \gamma_{\text{critical}}$.
\end{proof}

\subsection{Numerical Stability and Convergence}

\begin{theorem}[Numerical Convergence]
Let $a_h(t)$ be the numerical solution computed with time step $h$. Under standard regularity conditions (Lipschitz continuity of the right-hand side), for a $p$-th order method:
\[
\|a(t) - a_h(t)\| \leq C h^p,
\]
where $C$ is a constant depending on the integration interval and problem parameters.
\end{theorem}

\begin{remark}
These simulations illustrate how even modest modifications to the dark energy sector can have significant impacts on the overall expansion dynamics of the Universe. The interplay between matter and a decaying $\Lambda$ leads to rich dynamical behavior that may help resolve some of the tensions within the standard cosmological model.
\end{remark}

\subsection{Physical Interpretation}

\begin{definition}[Late-Time Behavior Classes]
We classify the late-time behavior of $a(t)$ into three classes based on $\gamma$:
\begin{enumerate}
\item \textbf{Accelerating} ($\gamma < \gamma_1$): $\lim_{t \to \infty} \ddot{a}(t) > 0$
\item \textbf{Coasting} ($\gamma = \gamma_1$): $\lim_{t \to \infty} \ddot{a}(t) = 0$
\item \textbf{Decelerating/Recollapsing} ($\gamma > \gamma_1$): $\lim_{t \to \infty} \ddot{a}(t) < 0$
\end{enumerate}
where $\gamma_1$ is a critical threshold.
\end{definition}

\begin{theorem}[Critical Decay Rate]
For the model with $\Omega_m = 0.3$ and $\Omega_{\Lambda} = 0.7$, there exists a critical decay rate $\gamma_{\text{crit}}$ such that:
\[
\gamma_{\text{crit}} = \frac{2}{3}H_0,
\]
above which the universe eventually recollapses.
\end{theorem}

\subsection{Observational Constraints}

\begin{proposition}[CMB Constraints]
Observations of the cosmic microwave background (CMB) constrain $\gamma$ to satisfy:
\[
\gamma < 10^{-2} H_0,
\]
based on the integrated Sachs-Wolfe effect and late-time structure formation.
\end{proposition}

\begin{remark}
Future observations from missions such as Euclid and the Vera Rubin Observatory will provide tighter constraints on dynamical dark energy models, potentially ruling out or confirming the exponential decay scenario.
\end{remark}
